% =====================================================================
%  The Boundary of Evaluator Auditing — NeurIPS Evaluations & Datasets (anonymous)
%  Batch 708 (paper 2 of the two-paper programme). Draft, not camera-ready.
%
%  WHAT THIS FILE IS
%    A *reorganisation* of results already produced in batches 697, 698 and 700
%    into a second, self-contained paper. NO new experiments were run for it:
%    every number traces to a frozen artifact of an earlier batch (see
%    data/708_paper2_number_audit.md for the number-by-number provenance table).
%
%  WHAT THIS FILE IS NOT
%    - It does NOT touch the first paper (queyi_neurips2027_v1.1.tex) or queyi_refs.bib.
%    - It runs no detect() call (detect_calls = 0).
%    - It claims no priority: 698-E found Manheim-Garrabrant (2018) for the
%      Goodhart taxonomy and arXiv:2608.24419 for construct-validity profiles.
%      Our phrasing throughout is "a subclass we name and formalise".
%
%  Template: official neurips_2025.sty, unmodified, used as a PLACEHOLDER
%  (the 2027 CFP is unpublished; the default notice is overridden below with a
%  neutral "under review" string). Target venue: TMLR (rolling) or NeurIPS 2027 E&D.
%
%  Page discipline (batch 708-C): the main text carries six tables; the
%  universality table, the T1 partition table, the full five-frame transition
%  matrix, the correctability-layer table, the theory-comparison table and all
%  supplementary curves live in the appendix.
% =====================================================================
\documentclass{article}

\usepackage[dandb]{neurips_2025}
\usepackage[utf8]{inputenc}
\usepackage[T1]{fontenc}
\usepackage{hyperref}
\usepackage{url}
\usepackage{booktabs}
\usepackage{amsfonts}
\usepackage{amsmath}
\usepackage{amssymb}
\usepackage{graphicx}
\usepackage{xcolor}
\usepackage{array}
\usepackage{amsthm}

\newtheorem{theorem}{Theorem}
\newtheorem{corollary}{Corollary}

\newcommand{\Rhat}{\hat R}
\newcommand{\Rstar}{R^{*}}
\newcommand{\calQ}{\mathcal{Q}}
\newcommand{\calG}{\mathcal{G}}
\newcommand{\catch}{\texttt{catch}}
\newcommand{\miss}{\texttt{miss}}
\newcommand{\unk}{\texttt{unknown}}
\newcommand{\T}{\mathsf{T}}
\newcommand{\pp}{\,\mathrm{pp}}

% Neutral notice string (the 2027 CFP is unpublished; do not print a wrong year).
\makeatletter
\renewcommand{\@noticestring}{%
  Under review (anonymized). Do not distribute.%
}
\makeatother

\title{The Boundary of Evaluator Auditing:\\
Identifiability, Correctability and Sample Complexity of Measurement Drift}

\author{%
  Anonymous Author(s) \\
  Affiliation \\
  Address \\
  \texttt{email} \\
}

\begin{document}
\maketitle

\begin{abstract}
Evaluation numbers are treated as facts about the systems they measure. We show that a large class of
evaluation failures requires \textbf{no optimisation pressure at all}: the number changes because the
\emph{caliber} of measurement changed, while the capability being measured does not. We call this
structural (passive) Goodhart and build an algebra for it.

We present a seven-axiom system for measurement drift and prove that only \textbf{two} of the axioms are
genuinely independent: idempotence, commutativity, non-invertibility and observability-stratification are
\emph{theorems} of the framework, and enumerability is a meta-requirement. We show the system is
\textbf{incomplete}: label-vocabulary drift (Type~III) cannot be expressed, and we register an eighth axiom
closing the label axis.

We then prove three boundary results. (i)~\textbf{Identifiability}: structural and real drift are not
decidable from paired data alone, yet are point-identified under nested calibers. (ii)~\textbf{Correctability}:
post-hoc repair is possible exactly when the drift acts on the post-measurement layer; acquisition-layer
drift requires re-measurement, with residual error up to $30.0\%$. (iii)~\textbf{Sample complexity}: the
paired design is \emph{powerless} against label and aggregation drift (discordance rate zero
$\Rightarrow$ infinite samples), while component-removal drift costs $849$ paired samples versus $17$ for
environment drift---cost is governed by the discordance rate, not the effect size.

We instantiate the framework on a C/C++ verification apparatus ($1{,}147$ samples) and, partially, on
LLM-as-a-judge evaluation (migration score $55.8/100$). All numbers come from frozen artifacts of earlier
batches; no new measurement was taken for this paper.
\end{abstract}

% =====================================================================
\section{Introduction}
\label{sec:intro}

Evaluation numbers are read as facts about the systems they measure. A detection rate, a judge score, a
benchmark accuracy: each is treated as a property of the artefact, portable across time and context. This
paper is about the cases where that reading fails for a reason unrelated to the artefact. The reported
number moves because the \emph{caliber} of the measurement moved---which components ran, which samples
entered the denominator, how per-sample verdicts were grouped, how they were aggregated---while the
capability under measurement is unchanged.

The classical literature on metric corruption assumes an optimiser. Goodhart's law
\citep{goodhart1975problems,strathern1997improving} is stated for a measure that \emph{is used to control} a
system; reward hacking \citep{skalse2022rewardhacking,amodei2016concrete} and reward-model over-optimisation
\citep{gao2023scaling} require a policy that maximises a proxy. \citet{manheim2018categorizing} catalogue
four variants, and the first is unambiguous: a metric \emph{which can be used to improve a system} is used
to such an extent that further optimisation is ineffective or harmful. \citet{majka2025goodhart} remove the
independence assumption and the learning paradigm but retain optimisation pressure; \citet{karwowski2023goodhart}
prove an optimal \emph{early stopping} rule---a rule that presupposes a trajectory to stop. We study the
subclass in which that premise is absent. Define \emph{structural (passive) Goodhart} as the conjunction of
three conditions: the reported value moved (SG-1), the mechanism-level capability did not (SG-2), and
\emph{no optimiser acts on the reported value} (SG-3). SG-3 separates this subclass from every formalisation
we could find, and it removes the object all classical remedies operate on: early stopping, regularisation,
robust optimisation, reward shaping and metric substitution all act on an optimisation trajectory, and when
there is no optimiser there is no trajectory.

That observation turns the subject into a measurement question, and measurement questions have boundaries.
This paper maps three of them.

\paragraph{Three boundaries.} \emph{Identifiability}: given paired measurements under two calibers, can we
tell whether the capability changed or only the caliber? Section~\ref{sec:ident} proves that in general we
cannot, and that under nested calibers we can. \emph{Correctability}: which drifts can be repaired after the
fact from retained records, and which require re-measurement? Section~\ref{sec:correct} proves a two-sided
criterion and quantifies the residual where repair fails. \emph{Sample complexity}: how many paired samples
does an audit need? Section~\ref{sec:sample} shows the answer is governed by how many samples \emph{carry}
the drift, not by how large the drift looks, and that for two of the four drift types the paired design has
exactly zero power at any sample size.

\paragraph{Contributions.}
\begin{itemize}\itemsep2pt
  \item \textbf{C1 (algebra and metatheory).} A seven-axiom algebra of measurement drift with six theorems
  (T1--T6). Its metatheory is settled by exhaustive model checking over all $2^{4\times3}=4096$ coverage
  structures $\times$ seven aggregation families: only A2 (monotonicity) and A5 (super-additivity) are
  genuinely independent; A3/A4/A6/A7 are theorems; A1 is a meta-axiom. The system is \emph{incomplete}---Type~III
  drift is inexpressible---so we register a missing A8 (label-axis closure) and extend the caliber signature
  to a 6-tuple. Caliber arbitrage is NP-hard (max-coverage reduction) with a greedy $(1-1/e)$ approximation.
  Sections~\ref{sec:algebra}--\ref{sec:meta}.
  \item \textbf{C2 (four types, and their universality classes).} Structural Goodhart comes in four types
  (denominator, environment, label, aggregation), each with its own control axis, and they do \emph{not}
  belong to one universality class: two are discontinuous, one is a smooth sub-critical decay with no
  critical point inside the observed range, and one is non-monotone. Section~\ref{sec:algebra}.
  \item \textbf{C3 (information-theoretic audit bound).} Paired designs have identically zero power against
  Type~III and Type~IV drift ($\psi=0$, hence $n=\infty$); Type~I costs $n=849$ and Type~II $n=17$ at
  $\alpha=0.05$ and power $0.8$. Audit cost is driven by the discordance rate, not the effect size.
  Section~\ref{sec:sample}.
  \item \textbf{C4 (validation, and a cross-domain arm).} Four experiments instantiate the theorems on a
  frozen C/C++ matrix ($1{,}147$ samples, $566$-sample paired split) and, partially, on an LLM-as-a-judge arm
  ($80$ samples $\times$ 2 judges $\times$ 2 prompts, temperature $0$), where the protocol migrates with score
  $55.84/100$---with the aggregation axis essentially non-transferable. Section~\ref{sec:exp}.
\end{itemize}

\paragraph{Scope and honesty.} The empirical basis is one C/C++ verification apparatus plus one partial LLM
arm; we claim no universality. We claim no priority: \citet{manheim2018categorizing} already taxonomise
Goodhart's law, and \citet{chen2026judge} already formalise construct-validity profiles for
LLM-as-a-judge---the closest neighbour, whose two-dimensional invariance/sensitivity profile answers a
different question (is the judge stable?) from ours (can the drift be repaired?). We claim no new
measurement: every number in Section~\ref{sec:exp} and the appendices was computed in an earlier batch from
a frozen matrix. Where a proof is incomplete we say so rather than fill the gap.

% =====================================================================
\section{Related Work}
\label{sec:related}

\paragraph{Goodhart's law and its variants.} The roots are \citet{goodhart1975problems} and
\citet{strathern1997improving}; \citet{muller2018tyranny} is the institutional critique. The formal line runs
through reward hacking \citep{skalse2022rewardhacking,amodei2016concrete}, over-optimisation scaling laws
\citep{gao2023scaling}, the four-variant taxonomy \citep{manheim2018categorizing}, the paradigm-agnostic
formalisation \citep{majka2025goodhart}, and the geometric treatment with an optimal early-stopping rule
\citep{karwowski2023goodhart}. All presuppose a system that optimises the measure, which is exactly what our
subclass removes (SG-3); consequently the classical remedies are undefined rather than merely ineffective.

\paragraph{Construct validity and measurement invariance.} Multitrait--multimethod matrices separate method
from trait variance \citep{campbell1959convergent}---the job our transition matrix does for caliber versus
mechanism; \citet{messick1995validity} argues validity is a property of inferences, not of scores; and
\citet{meredith1993measurement} formalises invariance across groups, with \citet{borsboom2006attack}
warning against defining the construct by the measurement model. The vocabulary has entered machine
learning \citep{bean2025constructvalidity,freiesleben2025benchmarkingepistemology}. The closest work is
\citet{chen2026judge}, who formalise a judge's construct validity as a two-dimensional profile---invariance
$S$ under construct-preserving edits and sensitivity $R$ under construct-changing edits---prove $S$ and $R$
independent, and show no scalar summary preserves all comparisons. We cite that approvingly and draw the
boundary explicitly: an $S/R$ profile is a \emph{probability profile} answering ``how stable is the judge?'';
our object is an algebra answering ``which caliber operations can be repaired post hoc, and which require
re-measurement?'' The latter question---our Theorem~\ref{thm:T6}---we did not find there. Statistical
concept drift \citep{cerqueira2026conceptdrift} concerns a changing data distribution, not a changing caliber.

\paragraph{Undecidability and detection limits.} That perfect static analysis is impossible is classical
\citep{turing1936computable,rice1953classes,landi1992undecidability,ramalingam1994undecidability}, with the
soundness/completeness trade-off \citep{cousot1977abstract} and the testing-versus-proof distinction
\citep{demillo1979social}. The precedent we follow is \citet{weyuker1988evaluation}, who uses axioms to
\emph{exclude} a perfect test-adequacy criterion; we do the same for a perfect drift detector (A2 excludes
minimum-type aggregations, A5 excludes mean-type aggregations, Theorem~\ref{thm:T6} excludes post-hoc
correction of acquisition-layer drift).

\paragraph{Evaluator auditing and meta-evaluation.} Existing work falls into five groups: comparisons of
evaluation \emph{methods} \citep{li2025meta,lv2026whoevaluates}; judge reliability audits
\citep{gu2024survey,tan2024judgebench,lambert2024rewardbench,zheng2023judging,norman2026reliabilityvalidity};
reporting templates \citep{mitchell2019modelcards,gebru2021datasheets}; adversarial-use audits; and
institutional audits \citep{costanzachock2022whoaudits,wang2026whoaudits,zhang2026whogrades}. Adjacent
negative results include audit failure modes \citep{li2026auditingaudit}, protocol-validity critiques
\citep{shao2026protocolvalidity}, sample-level auditing \citep{siedler2026samplelevel}, item-level data
release \citep{jiang2026itemlevel}, leaderboard dynamics \citep{shankar2025leaderboard}, audit-then-score
pipelines \citep{huang2026deepfact}, and reproducibility surveys
\citep{siddiq2025reproducibilitycrisis,antunes2024reproducibilitysurvey,pineau2021improving,eve2026validationcrisis}.
All five hold the caliber fixed: they audit results, not the act of measuring. We occupy a sixth
position---algebraising caliber drift---and claim no superiority over any of them
(\citealp{liang2023helm} audits models across a far broader scenario space than our single apparatus).

\paragraph{Software-verification and benchmark evaluation.} The apparatus sits in a mature tradition of
judging verifiers and test criteria \citep{klein2009sel4,leroy2009compcert,jia2011mutation,just2014mutants,
cadar2008klee,cui2024falsenegatives,li2024ubfuzz,beyer2026svcomp}, alongside a parallel tradition of
deliberate benchmark renewal \citep{kiela2021dynabench,white2024livebench,wang2024benchmarkevolving}. Both
treat the measurement configuration as fixed within a comparison; our point is orthogonal---it is the
configuration itself that is the variable.

\paragraph{Relation to our first paper.} A companion case study \citep{queyi2027caliber} audits one C/C++
verification apparatus and reports five findings, each an instance of one of the four drift types below. We
cite it for those findings only: that paper asks \emph{what was found}; this one asks what an audit can find
\emph{in principle}.

% =====================================================================
\section{Measurement Drift Algebra}
\label{sec:algebra}

\subsection{Apparatus and caliber}
\label{sec:apparatus}

A \emph{measurement apparatus} is a tuple $M=(D,A,E,\Theta,P)$, where $D$ is the sample population, $A$ the
asset pool (the components that acquire evidence), $E$ the environment, $\Theta$ the configuration, and $P$
the protocol (accounting rule, aggregation rule, label vocabulary). We call $Q(M)=(A,E,\Theta,P)$ the
\emph{caliber}; $D$ is not part of it. Two quantities must be kept apart. The \emph{reported rate} is
\begin{equation}
\Rhat(M)=\frac{\bigl|\{x: V(x;M)=\catch\}\bigr|}{\bigl|\{x: V(x;M)\in\{\catch,\miss\}\}\bigr|},
\label{eq:rhat}
\end{equation}
where $V(x;M)\in\{\catch,\miss,\unk\}$ is the per-sample verdict and $\unk$ is kept out of the denominator
under the unaware accounting rule; separating ``tested and missed'' from ``not tested'' is the
three-valued reading of the four-valued discipline used by the companion paper
\citep{belnap1977useful}. The \emph{mechanism-level capability} $\Rstar(M)$ is the rate the
apparatus would report under the reference caliber. $\Rstar$ is latent: it is what is being measured, not
what is reported. Every difficulty in Section~\ref{sec:ident} traces to that one fact.

Throughout, the instantiation is a C/C++ verification apparatus with $|A|=8$ declared assets (sanitizers,
compiler warnings, cross-compilation, linker, plus two components that are structurally $\unk$),
$|X|=1147$ samples, and coverage function $f(S)=|\bigcup_{a\in S}C_a|$, where $C_a$ is the set of samples
asset $a$ catches; the reference rate is $\Rhat(S)=f(S)/n$.

\subsection{Drift generators and the four types}
\label{sec:fourtypes}

A \emph{drift generator} is a map $\tau:M\to M'$ that changes the caliber but not $D$ and not the reference
apparatus; the \emph{drift} is $D_\tau(M)=\Rhat(\tau M)-\Rhat(M)$. Generators are classified by which
caliber component they act on.

\begin{table}[t]
\centering
\small
\caption{The four drift types and the measured drift on the frozen C/C++ matrix. Frames differ (Type~II is
a $566$-sample paired split; I, III, IV are computed on the full $1{,}147$-sample matrix) and are flagged
wherever the distinction matters.}
\label{tab:fourtypes}
\begin{tabular}{@{}c p{2.4cm} p{4.0cm} p{3.2cm} r@{}}
\toprule
Type & Acts on & Criterion & Instance & Drift \\
\midrule
I & denominator ($P$) & numerator fixed, denominator moves & accounting rule / degenerate slots & $+75.265\pp$ \\
II & environment ($E$) & assets gated by the environment & WSL $\to$ native ($n{=}566$) & $-35.3357\pp$ \\
III & label axis ($\lambda$) & per-sample verdicts fixed, grouping moves & $34$- vs $70$-type vocabulary & $-9.67\pp$ \\
IV & aggregation ($P$) & verdicts fixed, merge rule moves & OR vs per-asset mean & $+43.9700\pp$ \\
\bottomrule
\end{tabular}
\end{table}

The types are not mutually exclusive; they superpose. Three of the four numbers deserve a word. For Type~I,
the same $140$ catches are reported as $24.735\%$ under the unaware denominator and as $100.0\%$ when only
capability-relevant negatives enter it: a $+75.265\pp$ move produced by \emph{no new measurement}. For
Type~III the number is vocabulary-dependent in a way that cannot be rescaled---the share of defect types
whose blind-spot rate exceeds $50\%$ is $13/34=38.24\%$ under the normalised vocabulary and
$18/70=25.71\%$ under the original mixed one. For Type~IV the identical $1147\times 8$ verdict matrix
yields $61.6392\%$ under OR and $17.6692\%$ under the per-asset mean ($13.2518\%$ if the two structurally
$\unk$ assets are counted in the mean's denominator). The practical reading: a single headline number
cannot be compared across calibers, and a drift invisible under one aggregation rule can be $33.33\%$ under
another (Appendix~\ref{app:extra}, Table~\ref{tab:manip}).

\subsection{Universality classes: the four types are not one phenomenon}
\label{sec:universality}

It is tempting to speak of ``the phase transition of structural Goodhart''. The data do not support a single
one. Treating caliber deviation as a control parameter and measurement distortion $\varphi=|\Rhat-\Rstar|$
as an order parameter, the four types fall into three different classes (Appendix~\ref{app:extra},
Table~\ref{tab:universality}). For Type~I the eight-point scan $d=0,\dots,7$ gives
$g_{\text{app}}(d)=5.470342+5.284164\,d-0.259967\,d^{2}$ ($R^{2}=0.999724$) over $5.30\to29.88\pp$; the
response function $\chi=5.284164-0.519934\,d$ decreases monotonically and its stationary point
$d^{*}=10.1632$ lies \emph{outside} the data range $[0,7]$, so over the observable range there is no
critical point and Type~I is a smooth sub-critical decay. The fitted exponent $\beta=2.022681$ is a
consequence of the quadratic fit and must be read as model-implied. Consistently, the critical-slowing
indicators a genuine transition would show are absent: the increment variance and the lag-one
autocorrelation do not rise with $d$ but fall. What \emph{is} usable is structural---the count of zero-yield
components, and a label-fragmentation index
$\mathrm{FI}=(\#\text{singleton}+\#\text{two-sample classes})/\#\text{classes}$, which stood at $0.285714$
before a vocabulary normalisation and $0.0$ after, i.e.\ it was already signalling before the event.

We therefore downgrade C2 to a \emph{descriptive framework}: with $6$--$8$ data points and a model-implied
exponent, ``phase transition'' is an analogy. The practical consequence is asymmetric: a discontinuous drift
has no advance warning, because the jump is complete before any signal appears, so only design-time
structural indicators can pre-empt it; the smooth type is the one type where monitoring can work.

\subsection{The axiom system}
\label{sec:axioms}

Let $A$ be the asset set, $X$ the sample population, $C_a$ the catch set of asset $a$, and
$h(G):=f(A)-f(A\setminus G)$ the \emph{damage} of removing asset set $G$ (removal-type generator $\tau_G$).

\begin{table}[t]
\centering
\small
\caption{The seven axioms, and the metatheory verdict of Section~\ref{sec:meta}. The axioms are stated in
full here; the independence witnesses are in Table~\ref{tab:independence}.}
\label{tab:axioms}
\begin{tabular}{@{}c l p{8.6cm}@{}}
\toprule
\# & Axiom & Formalisation, and verdict \\
\midrule
A1 & caliber enumerability & each $\tau$ declares its type from a finite taxonomy; \emph{meta-axiom} \\
A2 & removal monotonicity & $G\subseteq H\Rightarrow\Rhat(\tau_H M)\le\Rhat(\tau_G M)\le\Rhat(M)$; \emph{genuine axiom} \\
A3 & idempotence & $\tau_G\circ\tau_G=\tau_G$; \emph{theorem} \\
A4 & commutativity, closure & $\tau_G\circ\tau_H=\tau_H\circ\tau_G=\tau_{G\cup H}$; \emph{theorem} \\
A5 & super-additivity of damage & $h(G\cup H)\ge h(G)+h(H)$, $G\cap H=\varnothing$; \emph{genuine axiom} \\
A6 & non-invertibility & the only invertible element is $\tau_\varnothing$; \emph{theorem} \\
A7 & observability stratification & silence is a property of $(\tau,\text{accounting})$, not of $\tau$; \emph{theorem} \\
\bottomrule
\end{tabular}
\end{table}

A3 and A4 give $(\{\tau_G\}_{G\subseteq A},\circ)\cong(2^{A},\cup)$: the generators form a commutative
idempotent monoid---a join-semilattice---and \emph{not} a group, because A6 says every non-identity element
lacks an inverse. That is the algebraic content behind the most practically useful corollary of the paper:
knowing $\Rhat(\tau_G M)$ and $G$ does not determine $\Rhat(M)$, since inversion would require the assets'
\emph{unique} catch sets---exactly what removal destroys. Drift cannot be corrected after the fact; it can
only be re-measured. Section~\ref{sec:correct} sharpens this into a two-sided criterion.

\subsection{T1: super-additivity with a closed-form residual}
\label{sec:T1}

\begin{theorem}[Super-additivity of removal-type damage]
\label{thm:T1}
Let $f$ be monotone and submodular, and let $S_0=A\cup B\cup L$ with $A,B,L$ pairwise disjoint. Write
$h(G)=f(S_0)-f(S_0\setminus G)$. Then $h(A\cup L)\ge h(A)+h(L)$, with closed-form residual
\begin{equation}
h(A\cup L)-h(A)-h(L)=\bigl|C_A\cap C_L\setminus C_B\bigr|\ \ge 0 .
\label{eq:t1residual}
\end{equation}
Equivalently, for the drift $D=-h/n$,
$D_{A\cup L}-(D_A+D_L)=-\bigl|C_A\cap C_L\setminus C_B\bigr|/n\le 0$.
\end{theorem}

The proof (Appendix~\ref{app:proofs}) expands
$m=|C_A\setminus C_B|+|C_L\setminus C_B|-|C_A\cap C_L\setminus C_B|$ against
$a=|C_A\setminus C_B|-|C_A\cap C_L\setminus C_B|$ and
$l=|C_L\setminus C_B|-|C_A\cap C_L\setminus C_B|$. Note the trap: the residual is \emph{not} the
intersection of the two assets' unique-contribution sets---those are disjoint by construction, so a test
built on them would pass for every input while testing nothing. On the $566$-sample split the closed form
matches the measured residual exactly in all three partitions tested, with one partition attaining
equality (Appendix~\ref{app:extra}, Table~\ref{tab:t1}), so the bound is tight here. The reading is
quantitative: when several calibers change in one evaluation their effects \emph{cannot be added up} to
correct for them---superposition \emph{understates} the total, by $4.95\pp$ in the partition where $28$
samples are catchable by both removals.

\subsection{T2--T4: non-invertibility, stratification, one-directionality}
\label{sec:T234}

\begin{theorem}[Non-invertibility]
\label{thm:T2}
For every $G\neq\varnothing$ there exist $M_1\neq M_2$ with $\tau_G M_1=\tau_G M_2$ (identical reported
verdict sequences). Hence $\tau_G$ has no inverse.
\end{theorem}

The construction takes $G=\{a\}$ and two apparatuses differing only in $a$'s catch set ($\varnothing$ versus
some non-empty $C_B$); after removal both verdict sequences are decided by $A\setminus\{a\}$, while their
reported rates differ. The generator is many-to-one.

\begin{theorem}[Silence is not intrinsic to the drift]
\label{thm:T3}
For a fixed $\tau$ there exist two accounting calibers with
$\Delta_{\unk}(\tau,\alpha_{\text{unaware}})=0$ but
$\Delta_{\unk}(\tau,\alpha_{\text{aware}})=-\mathrm{lost}/n\neq0$.
\end{theorem}

Under the unaware rule a sample with no catch is written $\miss$ even when the environment gate is
non-empty, so every lost catch lands in $\miss$ and the $\unk$ count does not move; under the aware rule the
same samples become $\unk$. Measured on the $566$-sample split: $0.0\pp$ under unaware and $+75.265\pp$
under aware. This corrects a hypothesis from an earlier design note in which $\Delta_{\unk}=0$ was proposed
as a sufficient-and-necessary fingerprint of silent degradation. By Theorem~\ref{thm:T3} it is a
fingerprint of the \emph{accounting rule}, not of the drift.

\begin{theorem}[One-directionality criterion]
\label{thm:T4}
If $\tau$ is purely removal-type ($S'\subseteq S$) then $T_{\miss\to\catch}=0$. The converse fails.
\end{theorem}

The forward direction is immediate: a catch under $S'$ implies a catch under the larger $S$. The converse
fails because a removal and an instrument upgrade can co-occur, giving both losses and gains with a net
$T_{\miss\to\catch}=0$. On the $566$-sample split all $11$ defect groups have
$c_g=T_{\miss\to\catch}=0$, with $b_g$ from $1$ (ODR/link) to $38$ (STL) and layerwise McNemar $p$ as small
as $7.3\times10^{-12}$---consistent with a purely removal-type environment drift. The falsification
condition is sharp and was already met once: a cross-tool comparison in the companion study produced $8$
reverse pairs and was correspondingly downgraded from a caliber-drift claim to a cross-regime stress test.

\subsection{T5: chains compose in the generator layer but not in the observation layer}
\label{sec:T5}

\begin{theorem}[Chapman--Kolmogorov condition for drift chains]
\label{thm:T5}
Let $E_1\supseteq E_2$ and let $T_{12},T_{23},T_{13}$ be per-sample three-state transition count matrices,
with $P_{23}$ the row-normalised $T_{23}$. Then
$T_{13}=T_{12}P_{23}\Longleftrightarrow T^{(E_2\to E_3)}_{\miss\to\catch}=0$, i.e.\ iff
$E_3\subseteq E_2$ (monotone absorption).
\end{theorem}

The proof is in Appendix~\ref{app:proofs}; the essential step is that monotone absorption makes
$P_{23}[\miss][\miss]=1$ and Theorem~\ref{thm:T4} kills the reverse pair, so the $(\catch,\miss)$ cell of
the product equals $T_{23}[\catch][\miss]+T_{12}[\catch][\miss]$, which is exactly the true $T_{13}$ cell by
screening. Measured on the $566$-sample split: the three nested chains have maximum absolute CK error
$0.000$, while the two non-nested chains (which restore a component) have $65.784$ and $74.272$, i.e.\
$11.6\%$ and $13.1\%$ of $n$. The error is mass-conserving---for the ``restore asan'' chain, CK moves
$65.784$ samples from $\catch{\to}\catch$ to $\catch{\to}\miss$ and conjures the same number of
$\miss{\to}\catch$ transitions that do not exist. The audit consequence is blunt: a generator chain always
composes (A4), but the \emph{observed} transition matrices do not, so an auditor cannot infer a multi-step
caliber change by multiplying matrices; each step must be re-measured in a paired design.

\subsection{T6: the correctability boundary (statement)}
\label{sec:T6statement}

\begin{theorem}[Correctability boundary]
\label{thm:T6}
Let $R=(V(x,a))_{x\in X,a\in A}\in\{\catch,\miss,\unk\}^{|X|\times|A|}$ be the original record and
$\rho=\mathrm{agg}(R)$ the report. A drift $\tau$ is repairable post hoc \emph{if and only if} either (1)
$\tau$ does not modify $R$ but only $\rho$ (accounting rule, aggregation rule, label vocabulary), so
retaining $R$ permits recomputing any $\rho$; or (2) $\tau$ deletes columns of $R$ that are determined
cellwise by the retained columns together with a prior. Otherwise $\tau$ is not repairable and re-measurement
is required.
\end{theorem}

Section~\ref{sec:correct} gives the layered verdict and the quantified residual error.

\subsection{Computational complexity}
\label{sec:complexity}

Once the caliber is enumerated, deciding whether a caliber change occurred is polynomial, $O(|T|\cdot n)$,
because silence is a two-cell equality test on the transition matrix (Theorem~\ref{thm:A3}). Deciding
whether the drift is structural or real is undecidable from paired data alone (Theorem~\ref{thm:A1}).
Choosing a caliber $S$ to push the reported value above a target, subject to a removal budget
$|A\setminus S|\le k$, is NP-hard by reduction from max-coverage \citep{feige1998threshold}---and the greedy
selection whose $(1-1/e)$ guarantee the companion study established is, read at the metatheoretic level,
precisely the standard greedy approximation for monotone submodular maximisation
\citep{nemhauser1978submodular,krause2014submodular}. The honest limit of that statement: it is a
reduction argument citing the classical hardness of max-coverage, not a formal reduction proved here, and at
the scale of the instantiation ($|A|=8$, $2^{8}=256$ calibers) the problem is trivially enumerable. Its
meaning is for evaluators with many components.

% =====================================================================
\section{Metatheory: which axioms are actually doing work}
\label{sec:meta}

\subsection{Exhaustive model checking}
\label{sec:metadomain}

Calling seven statements ``axioms'' invites the question of which are independent. We answered it by
exhaustive enumeration on a domain small enough to be complete: $|X|=4$ samples, $|A|=3$ assets, so
$2^{4\times3}=4096$ coverage structures, crossed with seven aggregation families deliberately chosen to
cover the monotone/non-monotone and submodular/non-submodular quadrants (union coverage; at-least-two; best
single; worst single; per-asset mean; per-asset sum; and exactly-one, which is non-monotone).

\begin{table}[t]
\centering
\small
\caption{Axiom independence over all $4096$ structures $\times$ $7$ aggregations. The minimal axiom set is
meta-axiom $\{A1\}$, structural axioms $\{A2,A5\}$, theorems $\{A3,A4,A6,A7\}$.}
\label{tab:independence}
\begin{tabular}{@{}c p{1.9cm} p{7.6cm}@{}}
\toprule
Axiom & Verdict & Counterexample family and witness \\
\midrule
A1 & meta-axiom & constrains the language, not the structure \\
A2 & genuine axiom & \texttt{min\_single}: $C_a{=}C_b{=}\{0\}, C_c{=}\varnothing$ gives $R(A){=}0$, but removing $c$ gives $1/4$ \\
A3 & theorem & none (always true): $\tau_G$ is column deletion, $A\setminus G\setminus G=A\setminus G$ \\
A4 & theorem & none (always true): set union commutes and is associative \\
A5 & genuine axiom & \texttt{at\_least\_2}: same structure, $h(\{a\}){=}h(\{b\}){=}1/4$ but $h(\{a,b\}){=}1/4<1/2$ \\
A6 & theorem & none (always true): column deletion is non-injective for every $G\neq\varnothing$ \\
A7 & theorem & none ($3840$ witnesses); provable once accounting is written into the structure (T3) \\
\bottomrule
\end{tabular}
\end{table}

This matters for a practical reason. Labelling a theorem as an axiom inflates the apparent strength of a
system and obscures where its genuine constraint lies. The real constraints here are that A2 requires the
report to be monotone in components---excluding minimum-type calibers---and that A5 requires the coverage
structure to be submodular---excluding at-least-two and mean-type calibers. Those two exclusions are the
axiom-level counterparts of the aggregation-drift manipulability results: a mean caliber violates A5
(adding a zero-yield component moves the report), and a minimum caliber violates A2 (removing the worst
component \emph{raises} the report).

We record an honest gap. A2 and A5 are mutually independent in the sense that A5's witness
(\texttt{at\_least\_2}) satisfies A2 while violating A5, so A5 does not imply A2; but within the seven-family
aggregation set we did \emph{not} find a witness violating A2 alone. A2's independence is therefore
established only against families that also violate A5, and a larger aggregation family would be needed for
a strict separation. The check as written is registered as incomplete.

\subsection{Incompleteness, and the missing A8}
\label{sec:A8}

The axiom system does not express Type~III drift, and we exhibit the witness rather than assert it. Fix the
coverage structure $C_a=\{0,1\}, C_b=\{2\}, C_c=\varnothing$ and the per-sample verdicts
$V=(\catch,\catch,\catch,\miss)$. Under a fine vocabulary the three classes have blind-spot rates
$\{0.0,0.0,1.0\}$; under a coarse vocabulary merging them, the single class has rate $\{0.25\}$. Every
per-sample verdict is unchanged, every quantity the seven axioms speak about is unchanged
($R_{\text{or}}=0.75$ throughout), yet the grouped statistics move. The reason is structural: the label
function $\lambda$ is \emph{not in the caliber signature}. The signature is $(A,E,\Theta,P)$ and has no slot
for it.

We therefore register the missing axiom and the extended signature.

\begin{quote}
\textbf{A8 (label-axis closure).} The caliber signature must be extended to $(D,A,E,\Theta,P,\lambda)$, and
every reported statistic must declare its grouping axis $\lambda$; statistics whose $\lambda$ is undeclared
are not comparable.
\end{quote}

This is not a stylistic requirement but the axiom form of a measured fact: the share of defect types
exceeding $50\%$ blind-spot rate is $38.24\%$ under the normalised $34$-type vocabulary and $25.71\%$ under
the original $70$-type one, and the relation is \emph{non-monotone} in granularity---the finer vocabulary
contains $14$ singleton classes, which inflate the denominator without inflating the numerator. There is no
direction in which one may ``correct'' a Type~III number from one vocabulary to another; it must be
recomputed.

\subsection{Comparison with other measurement theories}
\label{sec:metacompare}

The drift algebra sits beside three older measurement theories (Appendix~\ref{app:meta},
Table~\ref{tab:metacompare}). Two borrowings are worth naming. First, our transition matrix does the same
job as a multitrait--multimethod matrix \citep{campbell1959convergent}: it separates the ``method''
(caliber) from the ``trait'' (mechanism). Second, A8 is the analogue of measurement invariance
\citep{meredith1993measurement}: both require cross-group comparability to be declared and checked. The
difference is that psychometrics has mature tests (Lord's $\chi^{2}$, Mantel--Haenszel) where we have only a
structural requirement---a direction in which we could borrow. The correspondence is an analogy and we
claim no physical one: in particular, A6's irreversibility is a consequence of our defining $\tau$ as column
deletion, and would not be a theorem if $\tau$ were allowed to be an invertible instrument transform.

% =====================================================================
\section{Identifiability}
\label{sec:ident}

\subsection{The three-state transition matrix}
\label{sec:transition}

For two calibers $M,M'$ define the count matrix
$T^{(M\to M')}_{ij}=|\{x: V(x;M)=i,\ V(x;M')=j\}|$ for $i,j\in\{\catch,\miss,\unk\}$. It is the complete
sufficient statistic for a paired design, and it is what an audit should publish. For the environment change
$E_1\to E_2$ under unaware accounting the $\unk$ row and column are identically zero in \emph{all} frames,
and the five frames reconcile with the previously published values cell by cell ($4/4$ PASS, an independent
recomputation rather than a transcription; Appendix~\ref{app:extra}, Table~\ref{tab:transition}). On the
$566$-sample split: $E_1$ gives $340/226/0$ and $E_2$ gives $140/426/0$ (catch/miss/unknown), so
$T_{\catch\to\catch}=140$, $T_{\catch\to\miss}=200$, $T_{\miss\to\catch}=0$, $T_{\miss\to\miss}=226$, and
McNemar $b=200$, $c=0$, $p=1.245\times10^{-60}$.

Row-normalising gives $P[\catch][\catch]=0.4118$, $P[\catch][\miss]=0.5882$ and $P[\miss][\miss]=1.0000$.
The reading motivates the whole paper: $58.8\%$ of the samples the reference caliber caught are reported as
\emph{tested and missed} under the reduced caliber, and $0\%$ are reported as \emph{not tested}. Lost
capability is disguised as a lower rate, not as a gap in coverage.

The paired design also exposes a window in which degradation is nearly invisible. Over the $63$ non-empty
removal subsets, $|W(\rho)|=\{\calG: \text{retention}\ge\rho\wedge\text{unsound negatives}>0\}$ is $0$ at
$\rho=100\%$ and $99\%$, $1$ at $95\%$, $7$ at $90\%$, $18$ at $80\%$ and $49$ at $50\%$. The widest window
is removing the \emph{linker} alone: retention $98.8236\%$, so the headline rate barely moves, yet $22$
negative verdicts are false ($14$ ODR/link, $8$ legacy) and type-conditional credibility falls to
$90.43\%$. Reporting a single rate cannot see this. A hypothesis that the window's severity is governed by
how \emph{concentrated} the lost capability is survives for the window's \emph{position}
($\rho(\text{concentration},\text{retention})=+0.9429$, $n=6$) but is refuted for its \emph{severity} across
all $63$ configurations ($\rho(\text{loss},\text{unsound})=+0.9950$ versus
$|\rho(\text{concentration},\text{unsound})|=0.6662$). We report both, and we report that the $n=6$
correlation is exploratory.

\subsection{Structural versus real drift is not decidable}
\label{sec:A1}

\begin{theorem}[Non-identifiability from paired data]
\label{thm:A1}
Given paired observations $(M,M')$ alone, it is not decidable whether the drift is structural
($\Rstar$ unchanged) or real ($\Rstar$ changed).
\end{theorem}

\begin{proof}[Proof sketch]
Let $\tau:q_2\to q_1$ be removal-type and suppose the observed table has $T_{\catch\to\miss}=b>0$,
$T_{\miss\to\catch}=c=0$. Construct two data-generating processes. $G_{\text{struct}}$ sets $c(x)=1$ iff
$V(x;q_2)=\catch$, so $\Rstar$ equals the reference rate and is unchanged; $G_{\text{real}}$ sets $c(x)=1$
iff $V(x;q_1)=\catch$, so $\Rstar$ falls. Both produce the same joint distribution on
$(V(\cdot;q_1),V(\cdot;q_2))$---the same $T$, the same $b$, the same $c=0$---and opposite conclusions about
the capability. The latent capability is unobserved and no constraint pins it to the observed
distribution.
\end{proof}

This is not pessimism; it is the standard conclusion for an unidentified latent quantity. It cannot be
broken by more statistics, only by another measurement.

\subsection{Nested calibers are point-identified; non-nested ones are not}
\label{sec:A2}

\begin{theorem}[Identifiability under nesting]
\label{thm:A2}
Let the measurement function be the coverage function $f(S)=|\bigcup_{a\in S}C_a|$ (monotone and submodular,
verified empirically in the companion study). (i) If $q_1\subseteq q_2$ then $\Rstar$ has been observed at
the reference caliber $q_2$, so the drift is point-identified:
$D_\tau=\Rhat(q_1)-\Rhat(q_2)\le0$ by monotonicity. (ii) If $q_1$ and $q_2$ are incomparable, the reference
caliber $q^{*}=q_1\cup q_2$ is unobserved and only a submodular interval is available,
$\max(\Rhat(q_1),\Rhat(q_2))\le\Rhat(q^{*})\le\Rhat(q_1)+\Rhat(q_2)-\Rhat(q_1\cap q_2)$; point
identification fails and the verdict is \emph{unresolved}.
\end{theorem}

\begin{corollary}
\label{cor:A2.1}
All four drift instances of Table~\ref{tab:fourtypes} are nested, hence their drift magnitudes are
mathematically point-identified. What is genuinely unidentified is the separation of composition from
mechanism effects within a single reported difference---a decomposition into two unobservable components,
which no additional observation of the caliber can fix.
\end{corollary}

Corollary~\ref{cor:A2.1} has a rhetorical consequence we take seriously: the companion study's inability to
separate composition from mechanism should be written as a formal limitation (non-identifiability), not as a
deferral.

\begin{theorem}[Decidability dichotomy]
\label{thm:A3}
(a) Whether $\tau$ is silent is decidable in $O(n)$: it is the equality test
$T_{\catch\to\unk}=T_{\miss\to\unk}=0$ on the transition matrix. (b) Whether the drift is structural is
undecidable from $(M,M')$ alone, unless the calibers are nested or the reference caliber is measured
(Theorems~\ref{thm:A1}, \ref{thm:A2}).
\end{theorem}

The two halves are of different difficulty and must not be conflated: a test that decides silence cheaply
says nothing about structure, and a system that is silent is not thereby degenerate.

\subsection{Per-asset sensitivity and the contaminated-negative rate}
\label{sec:sensitivity}

Which asset's removal manufactures false negatives? Not the highest-yield one but the one with the largest
\emph{unique} catch set, since samples also caught by a surviving asset do not change the OR verdict. On the
$1{,}147$-sample matrix: asan $139$ unique catches ($12.12\%$), tsan $85$ ($7.41\%$), ubsan $83$ ($7.24\%$),
compiler-warn $53$, cross-compile $14$, linker $10$, and the two structurally $\unk$ assets $0$. The
contaminated-negative rate---reported negatives that are in fact positives---is $46.95\%$ on the
$566$-sample split ($200$ of $426$) and $50.89\%$ on the $1{,}147$-sample matrix ($456$ of $896$). The two
numbers are both correct and \emph{not interchangeable}: they come from different frames, and even the
asset ranking differs between them (ubsan $>$ tsan on $566$; tsan $>$ ubsan on $1147$)---itself a
composition-drift instance acting on an environment-drift conclusion. Any citation of these numbers must
carry the frame.

% =====================================================================
\section{Correctability}
\label{sec:correct}

Theorem~\ref{thm:T2} said removal-type drift is irreversible. Theorem~\ref{thm:T6} sharpens that into an
operational dichotomy: irreversibility is not a property of all removal, but of removing columns that carry
information the retained columns cannot supply. Post-processing drifts---accounting rule, aggregation rule,
label vocabulary---leave the record $R$ intact and are repaired at zero residual cost; so are the
structurally-$\unk$ columns, whose values follow from a prior. Environment-gated removal of the sanitizer
columns is in the other class.

To quantify the failure we gave post-hoc repair its best chance: predict the removed column from the verdict
combination of the retained columns (cw, cc, linker) by lookup on the $566$-sample split.

\begin{table}[t]
\centering
\small
\caption{Best post-hoc reconstruction of removed sanitizer columns. ``Baseline'' is the trivial rule that
always guesses $\unk$. Even with the best lookup, asan retains $30.04\%$ unrecoverable verdicts.}
\label{tab:residual}
\begin{tabular}{@{}lrrrr@{}}
\toprule
Removed column & best accuracy & \textbf{residual error} & baseline & gain \\
\midrule
asan & $69.96\%$ & $\mathbf{30.04\%}$ & $0.71\%$ & $+69.26\pp$ \\
ubsan & $78.62\%$ & $\mathbf{21.38\%}$ & $0.71\%$ & $+77.92\pp$ \\
tsan & $78.27\%$ & $\mathbf{21.73\%}$ & $0.71\%$ & $+77.56\pp$ \\
wunsequenced & $100\%$ & $0.00\%$ & $100\%$ & $0$ (prior) \\
compile-time & $100\%$ & $0.00\%$ & $100\%$ & $0$ (prior) \\
\bottomrule
\end{tabular}
\end{table}

The contrast within Table~\ref{tab:residual} is the content of the theorem. Two removed columns are
recovered perfectly at zero cost because their values are determined by a prior (they are structurally
$\unk$). Three are not, because they carry unique information: about a third of asan's verdicts, and about a
fifth of ubsan's and tsan's, cannot be recovered from anything retained. ``Not correctable'' is therefore
not a property of removal as such but of removing a column that carries information nothing else carries---
and the criterion is computable \emph{before} the removal is made. The operational rule that follows is a
design constraint rather than a diagnostic: retain the per-asset $\times$ per-sample verdict matrix, because
it is the only record from which post-processing drifts are repairable.

% =====================================================================
\section{Sample Complexity}
\label{sec:sample}

\subsection{Formalisation}
\label{sec:sampleform}

Test $H_0$: no drift (equal marginal rates) against $H_1$: drift with marginal difference at least
$\varepsilon$, using a paired design and a sign test on the discordant pairs (equivalently, an exact
McNemar test \citep{mcnemar1947note}). Let $\psi=\Pr[\text{the two calibers disagree}]$ be the discordance
rate and $\pi=\Pr[\text{direction}\mid\text{disagreement}]$. The identity
\begin{equation}
\varepsilon=\psi\,(2\pi-1),
\label{eq:epspsi}
\end{equation}
is the hinge of the analysis. With $\alpha=0.05$ and power $0.8$
\citep{cohen1988statistical} ($z_{1-\alpha/2}=1.9600$, $z_{1-\beta}=0.8416$), the required number of
discordant pairs and samples are
\begin{equation}
m\ \ge\ \frac{\bigl(z_{1-\alpha/2}\sqrt{1/4}+z_{1-\beta}\sqrt{\pi(1-\pi)}\bigr)^{2}}{(\pi-1/2)^{2}},
\qquad n=\frac{m}{\psi},
\label{eq:sample}
\end{equation}
the paired-design form of the standard sample-size calculation \citep{connor1987sample}, with an
information-theoretic floor $m\ge\log(1/\beta)/\mathrm{KL}(\pi\|1/2)$. For the fully
one-directional case $\pi=1$---not a degenerate corner but the actual situation on this data, where
$c=0$---an exact binomial criterion replaces the normal approximation:
$m\ge\lceil\log(\alpha/2)/\log(1/2)\rceil=6$.

\subsection{The four types cost wildly different amounts}
\label{sec:sampleresults}

\begin{table}[t]
\centering
\small
\caption{Sample complexity by drift type on the $566$-sample paired split, $\alpha=0.05$, power $0.8$.
Types~III and IV have $\psi=0$: their caliber change leaves \emph{every} per-sample verdict unchanged, so
no paired test has power at any sample size.}
\label{tab:sample}
\begin{tabular}{@{}c p{4.6cm} r r r r r@{}}
\toprule
Type & Caliber change & $\psi$ & $\varepsilon$ & $m$ & $\mathbf{n}$ & KL floor \\
\midrule
I & drop 3 degenerate assets & $0.0071$ & $0.0071$ & $6$ & $\mathbf{849}$ & $329$ \\
II & environment $E_1\to E_2$ & $0.3534$ & $0.3534$ & $6$ & $\mathbf{17}$ & $7$ \\
III & label vocabulary change & $0.0000$ & $0.0000$ & --- & $\boldsymbol{\infty}$ & --- \\
IV & aggregation rule change & $0.0000$ & $0.0000$ & --- & $\boldsymbol{\infty}$ & --- \\
\bottomrule
\end{tabular}
\end{table}

Three consequences follow, each a correction to a natural expectation.

\paragraph{(i) Types III and IV are not hard---they are undetectable by this design.} Their $\psi$ is zero
because the caliber change moves the report without moving any per-sample verdict
(Section~\ref{sec:A8} for III; Section~\ref{sec:fourtypes} for IV). Every test built on per-sample
verdicts---paired McNemar, permutation, bootstrap---has power identically zero. This is a property of the
design, not of the sample size, and no budget fixes it.

\paragraph{(ii) Cost is governed by the discordance rate, not the effect size.} Types I and II have the
\emph{same} $\varepsilon$, both being fully one-directional; but $n=m/\psi$, and Type~I's $\psi$ is $1/50$
of Type~II's, so its sample requirement is $50$ times larger ($849$ versus $17$). The mechanism is
concrete: Type~I's drift is carried by only $4$ unique catches of the linker in a $566$-sample frame, so $6$
discordant pairs require measuring $849$ samples. The practical rule is that audit budgets should be planned
by \emph{how many samples carry the drift}, not by how large the drift looks: a drift reading $-35\pp$ can
be cheap, and one reading $0.71\pp$ can be expensive.

\paragraph{(iii) The paired step must become a dual design.} A five-step audit protocol in our earlier work
treats sample-level pairing as the universal step. The result above shows it covers only the $A$ and $E$
axes. The label and aggregation axes need a \emph{design-level contrast}: one set of per-sample verdicts,
two label functions (or two aggregation rules), and a comparison of \emph{grouped statistics} rather than
per-sample verdicts. The corrected protocol therefore has two designs, with cost model $n=m/\psi$ estimated
per axis.

\subsection{Sampling strategy}
\label{sec:strategy}

With $400$ repetitions over eight budgets and three strategies on the $566$-sample frame, Neyman allocation
using true stratum discordance rates dominates at moderate budgets ($0.863$ versus $0.777$ for uniform at
$B=20$); stratified-proportional sits between; all strategies saturate by $B\ge50$. Neyman with true rates
is an upper bound, not an implementable policy, and an earlier two-stage adaptive variant that estimates the
rates from a pilot quarter is \emph{worse} than uniform at small budgets ($0.280$ versus $0.505$ at
$B=15$), because the effect here is large and the pilot's fixed cost is not repaid. The gap between the best
simulated strategy and the theoretical requirement of Table~\ref{tab:sample} is due to the discreteness of
the exact McNemar test at small budgets: reaching theoretical power needs a budget two to three times the
theoretical $n$.

% =====================================================================
\section{Empirical Validation}
\label{sec:exp}

Every number in this section was computed in an earlier batch from a frozen matrix; no new measurement was
taken for this paper. The experiments test the theorems, and each names the statement it checks.

\paragraph{E1: the algebra on a frozen matrix.} Setting: the C/C++ apparatus with $|A|=8$ and frames
$1147$ (capability boundary), $1137$ (all-pool), $566$ (paired evaluation split) and $110$ (real-world).
Method: recompute the transition matrices from the frozen verdicts and evaluate T1 on three partitions.
Result: the five transition matrices are reproduced with $\unk$ identically zero and reconcile with the
previously published values in all four comparable frames (nine quantities per frame, all equal); T1's
closed form matches the measured residual exactly in all three partitions, one of which attains equality;
T4 holds layerwise (all $11$ defect groups have $c_g=0$); and T3's predicted accounting dependence
reproduces ($0.0\pp$ unaware versus $+75.265\pp$ aware). This tests T1, T3, T4 and the identifiability
dichotomy of Theorem~\ref{thm:A3}.

\paragraph{E2: the correctability boundary.} Setting: the $566$-sample paired split, five drift chains
(three nested, two non-nested that restore a component). Method: compare the composed transition matrix
against the directly measured one, then reconstruct each removed column from the retained columns by
lookup. Result: nested chains satisfy Chapman--Kolmogorov exactly (maximum absolute error $0.000$, $3/3$);
non-nested chains violate it by $65.784$ and $74.272$ ($11.6\%$ and $13.1\%$ of $n$) with the error
mass-conserving; best post-hoc reconstruction leaves residual errors of $30.04\%$ (asan), $21.38\%$ (ubsan)
and $21.73\%$ (tsan) against $0.71\%$ for the trivial baseline, while the two structurally-$\unk$ columns
are recovered perfectly. This tests T5 and T6 jointly.

\paragraph{E3: sample complexity.} Setting: the $566$-sample frame, three sampling strategies $\times$
eight budgets $\times$ $400$ repetitions, exact McNemar at $\alpha=0.05$. Method: empirical detection
probability versus the closed-form requirement of Eq.~\eqref{eq:sample}. Result: Types~III and IV have
$\psi=0$ and $n=\infty$---the paired design is powerless, not merely underpowered; Type~I requires $n=849$
(KL floor $329$) and Type~II $n=17$ (KL floor $7$). This tests all of Section~\ref{sec:sample}, and it is
the experiment that changes the audit protocol.

\paragraph{E4: a cross-domain arm (partial).} Setting: an LLM-as-a-judge arm with $80$ samples $\times$
$2$ judges (a reasoning model and a non-reasoning model) $\times$ $2$ semantically equivalent prompts,
temperature $0$, with the declared per-sample verdict label as ground truth. This is a partial
instantiation: the planned cross-domain data collection was not completed, so the arm reuses an
in-repository judge study rather than public judge benchmarks.

\begin{table}[t]
\centering
\small
\caption{Cross-domain arm. The aggregation axis is essentially non-transferable: a change of aggregation
rule moves the LLM report by $1.03\pp$ against $33.33\pp$ in the C/C++ apparatus. Components are scored out
of $25$ with self-set weights and no external calibration.}
\label{tab:crossdomain}
\begin{tabular}{@{}p{2.7cm} p{4.0cm} p{3.0cm} r@{}}
\toprule
Axis & C/C++ apparatus & LLM arm & Migration \\
\midrule
caliber (prompt) & $-35.34\pp$ & invariance $87.5\%$ / $96.25\%$ & $15.62/25$ \\
environment (judge model) & two environments & false-report gap $13.0\pp$ & $16.25/25$ \\
aggregation & $33.33\pp$ & $1.03\pp$ & $2.58/25$ \\
label & $38.24\%$ vs $25.71\%$ & swing $8.56\pp$ & $21.39/25$ \\
\midrule
\textbf{total} & & & $\mathbf{55.84/100}$ \\
\bottomrule
\end{tabular}
\end{table}

The structural-Goodhart rate in the LLM arm---the fraction of samples a judge calls a catch while the
declared label says miss---is $26.25\%$ (reasoning judge) and $42.50\%$ (non-reasoning judge). The ten most
severe cases concentrate in the \emph{embedded} (6) and \emph{concurrency} (3) defect families, exactly the
families requiring an external specification. That agreement with the companion study's ``needs-oracle''
classes is suggestive rather than conclusive: human inter-annotator agreement in this arm is zero, so
``judge is wrong'' and ``label is wrong'' cannot be separated, and the severe cases must not be read as an
error list.

The arm also produced a live instance of Theorem~\ref{thm:T6}. The raw judge log lacks a run-identifier
field, and one model was run twice at different token budgets, so runs can only be inferred from completion
length. That inference moves the cross-model agreement rate by up to $39.19\pp$; the axis is therefore not
reliably estimable, and we fall back to published aggregate quantities. The missing metadata field is an
acquisition-layer defect, and by Theorem~\ref{thm:T6} it is not repairable after the fact---the arm
illustrates the theorem on itself. We record two further limits: the arm is not a full cross-domain
validation, and it records no model snapshot identifier, so its conclusion has a short reproduction horizon
---itself a limitation of the theorem's scope.

% =====================================================================
\section{Discussion}
\label{sec:discussion}

\paragraph{The three boundaries decide what auditing can do.} Identifiability says that from paired data
alone an auditor cannot separate a changed capability from a changed caliber, and that nesting or a
reference measurement is what breaks the tie. Correctability says that post-processing drift can be repaired
from retained records while acquisition-layer drift cannot, with a computable residual where repair fails.
Sample complexity says the cost of an audit is set by how many samples carry the drift, and that for two of
the four types the paired design is powerless by construction. Together these bound an audit from both
sides: what it can establish, and what it must buy.

\paragraph{Relation to the companion paper.} The companion study reports what was found in one apparatus;
this paper reports what an audit can find in principle. Its five findings are each an instance of one of the
four drift types: the collapse of an apparent selection gain is composition drift; the environment-induced
fall in real-defect detection is environment drift; the vocabulary-dependent blind-spot share is label
drift; and the headline blind-spot percentage is a consequence of the aggregation caliber. We cite the
companion for those findings and do not restate its arguments.

\paragraph{Limitations.} The empirical basis is one domain (C/C++ verification) plus one partial LLM arm;
the environment axis has two control points; $92.9\%$ of the corpus is human-planted; the ``always true''
metatheory verdicts hold only over the $4096$-structure, seven-aggregation domain, so they are in-domain
evidence rather than general proofs; A2 and A5 could not be separated within that domain; the Type~I
exponent $\beta$ is model-implied; the phase-transition language is an analogy with no thermodynamic
counterpart; the LLM arm has two judges, two prompts and eighty samples with zero human agreement; and the
risk thresholds and migration-score weights are self-set with no external calibration.

\paragraph{Future work.} Four items follow directly. (1) Formalise A8 and test the extended signature on a
second domain. (2) Derive the sample complexity of \emph{non-paired} designs, which is what Types~III and IV
actually require. (3) Add a third environment profile (a container image) to lift the environment axis from
two control points to three. (4) Build an equi-marginal replacement design to separate composition from
mechanism effects---which would convert Corollary~\ref{cor:A2.1}'s non-identifiability into identifiability
using a read-only rearrangement rather than a new measurement.

\paragraph{What an evaluator should do on Monday.} The three cheapest interventions are also the most
effective. Remove zero-yield components from the \emph{declared} asset pool, keeping them as a capability
boundary record: this eliminates composition inflation entirely, because the removed components catch
nothing ($0$ of $1147$). Switch the accounting rule from unaware to aware: this makes the contaminated
negatives visible rather than silent. Declare the aggregation rule and attach a manipulability check: this
turns aggregation drift from invisible into measurable. All three cost approximately nothing---a declaration
change, an already-implemented accounting mode, and an already-implemented scorer---and together they raise a
measured anti-drift score from $53.05$ to $66.67$ out of $100$ in a simulation over the same frozen matrix.
In the same simulation, a design that reduces the asset count for simplicity scores $0.3389$ against
$0.6230$ for the status quo and $0.7721$ for the theory-implied design, and its composition inflation is
$1.0000$: with the pool reduced to the selection budget there is no selection problem left to measure.

% =====================================================================
\section{Conclusion}
\label{sec:conclusion}

Evaluation numbers can change without any optimisation pressure, because the caliber of measurement changed
while the capability did not. We gave that structure an algebra---seven axioms of which only two are
independent---and proved its three boundaries: structural and real drift are not identifiable from paired
data alone, post-hoc repair is possible exactly on the post-measurement layer, and audit cost is governed by
the discordance rate rather than the effect size, with two of the four drift types undetectable by a paired
design at any sample size. The most useful practical consequence is that the three cheapest
interventions---declaring the pool without zero-yield components, accounting for untested samples, and
declaring the aggregation rule---cost approximately nothing and raise the measured anti-drift score from
$53.05$ to $66.67$.

\bibliographystyle{plainnat}
\bibliography{paper2_refs}

\appendix

% =====================================================================
\section{Proofs}
\label{app:proofs}

\subsection*{Theorem~\ref{thm:T1} (super-additivity, closed-form residual)}

Let $C_S:=\bigcup_{a\in S}C_a$ and write
$a_0:=|C_A\setminus(C_B\cup C_L)|$, $l_0:=|C_L\setminus(C_A\cup C_B)|$,
$m_0:=|(C_A\cup C_L)\setminus C_B|$. Then $h(A)=a_0$, $h(L)=l_0$, $h(A\cup L)=m_0$, and
\begin{equation}
m_0=|C_A\setminus C_B|+|C_L\setminus C_B|-|C_A\cap C_L\setminus C_B|,
\end{equation}
while
\begin{equation}
a_0=|C_A\setminus C_B|-|C_A\cap C_L\setminus C_B|,\qquad
l_0=|C_L\setminus C_B|-|C_A\cap C_L\setminus C_B|,
\end{equation}
because the part of $C_A\setminus C_B$ that lies in $C_L$ is exactly $C_A\cap C_L\setminus C_B$, and
symmetrically for $C_L$. Subtracting gives
$m_0-(a_0+l_0)=|C_A\cap C_L\setminus C_B|\ge0$, which is Eq.~\eqref{eq:t1residual}.

Equivalently, in the language of submodularity: $f$ submodular implies $g(S):=f(S_0\setminus S)$ submodular
(substitute $X=S_0\setminus S$, $Y=S_0\setminus T$ into the submodular inequality), so $h=f(S_0)-g$ is
supermodular and $h(A\cup L)\ge h(A)+h(L)-h(A\cap L)=h(A)+h(L)$ since $A\cap L=\varnothing$. $\square$

\paragraph{A trap worth recording.} The residual is \emph{not} the intersection of $A$'s unique-contribution
set with $L$'s. Those two sets are disjoint by construction---one excludes $C_L$ entirely, the other
contains only elements of $C_L$---so their intersection is identically empty and a test built on it would
pass for every input while testing nothing. The correct form is $|C_A\cap C_L\setminus C_B|$, which can be
non-zero precisely when a third asset set $B$ is present.

\subsection*{Theorem~\ref{thm:T2} (non-invertibility)}

Take $G=\{a\}$ for any asset $a$. Build $M_1,M_2$ agreeing with $M$ on all assets except $a$, with
$C_a^{(1)}=\varnothing$ and $C_a^{(2)}=C_B$ for some non-empty retained asset set $B$. After removing $a$,
both OR verdict sequences are decided by $A\setminus\{a\}$ and are therefore identical, while
$\Rhat(M_1)<\Rhat(M_2)$. So $\tau_G$ is many-to-one and has no inverse. $\square$

\subsection*{Theorem~\ref{thm:T3} (stratification of silence)}

Under unaware accounting, a sample with no catch is written $\miss$ even when the environment gate is
non-empty, so every catch lost to removal lands in $\miss$ and the $\unk$ count is unchanged:
$\Delta_{\unk}=0$. Under aware accounting the same samples become $\unk$, so
$\Delta_{\unk}=+\mathrm{lost}/n\neq0$. Both rules are admissible; neither is ``truer''. Aware makes the
problem visible at the cost of driving the conditional-recall denominator towards zero; unaware makes it
invisible at the cost of a report that misstates. $\square$

\subsection*{Theorem~\ref{thm:T5} (Chapman--Kolmogorov condition)}

($\Rightarrow$) Let $E_3\subseteq E_2\subseteq E_1$ and $V_i=V(x;E_i)$. \emph{Absorption}: if $V_2=\miss$
then no $a\in E_2$ catches $x$; since $E_3\subseteq E_2$, no $a\in E_3$ catches $x$ either, so
$V_3\in\{\miss,\unk\}$ (here $\unk\equiv0$), giving $P_{23}[\miss][\miss]=1$. \emph{One-directionality}:
by Theorem~\ref{thm:T4}, $T_{12}[\miss][\catch]=0$. The $(\catch,\miss)$ cell of the product is then
$T_{12}[c][c]\cdot P_{23}[c][m]+T_{12}[c][m]\cdot1$, and
$P_{23}[c][m]=T_{23}[c][m]/|E_2\text{-catch}|$ with
$|E_2\text{-catch}|=T_{12}[c][c]+T_{12}[m][c]=T_{12}[c][c]$, so the first term is $T_{23}[c][m]$ and the
product cell equals $T_{23}[c][m]+T_{12}[c][m]$. The true cell is
$\#\{V_1{=}c,V_3{=}m\}=\#\{V_2{=}c,V_3{=}m\}+\#\{V_2{=}m,V_3{=}m\}=T_{23}[c][m]+T_{12}[c][m]$ by
absorption. The other cells follow from the row sums and total conservation.

($\Leftarrow$) If $T_{23}[m][c]>0$ there is a sample with $V_2=\miss$, $V_3=\catch$. If
$E_3\not\subseteq E_2$, $E_3$ contains an asset absent from $E_2$, so this sample's $V_1$ may be either
$\catch$ or $\miss$: $V_2$ no longer screens $V_1$, the screening condition fails, and the product's
$(\miss,\catch)$ cell is non-zero whereas the true $T_{13}$ cell is $0$. $\square$

\subsection*{Theorem~\ref{thm:T6} (correctability boundary)}

(1) If $\tau$ acts only on $\rho$ via a known map $\rho\mapsto\rho'$ and $R$ is retained, then $\rho'$ is
computable, and $\rho$ is recoverable either directly (if $\rho\mapsto\rho'$ is injective on its range) or
from $R$; residual error $0$. (2) If the deleted columns are cellwise determined by the retained columns
together with a prior---as the structurally-$\unk$ columns are---they are reconstructible cell by cell;
residual error $0$. (3) Otherwise the deleted columns carry information not determined by the retained ones,
so by the construction of Theorem~\ref{thm:T2} there exist $R_1\neq R_2$ with $R_1'=R_2'$ after deletion:
the map is many-to-one, has no inverse, and no post-hoc procedure recovers $\rho$ in general. The residual
error is $1$ minus the best achievable prediction of the deleted column from the retained ones, which
Table~\ref{tab:residual} evaluates. $\square$

% =====================================================================
\section{Metatheory: full tables}
\label{app:meta}

Model domain: $|X|=4$ samples, $|A|=3$ assets, all $2^{4\times3}=4096$ coverage structures, seven
aggregation families---\texttt{or\_union} (monotone, submodular), \texttt{at\_least\_2} (monotone,
non-submodular), \texttt{max\_single}, \texttt{min\_single}, \texttt{mean\_single} (all monotone,
non-submodular), \texttt{sum\_single} (monotone, additive), and
\texttt{exactly\_one\_nonmonotone} (non-monotone, the source of A2 counterexamples). Drift operators are
removal-type, $\tau_G$ retaining $A\setminus G$, and damage is $h(G)=R(A)-R(A\setminus G)$.

\begin{table}[h]
\centering
\small
\caption{Completeness check. The label-coarsening witness shows a model in which every axiom-relevant
quantity is fixed while the grouped statistics move---so Type~III drift is inexpressible in the current
signature.}
\label{tab:completeness}
\begin{tabular}{@{}p{5.4cm} p{1.6cm} p{4.4cm}@{}}
\toprule
Check & Result & Note \\
\midrule
all four observed drift types expressible? & no & Type~III cannot be expressed \\
model satisfying all axioms yet drifting? & yes & label-coarsening witness \\
is a ninth axiom needed? & not found beyond A8 & not excluded: the domain is finite \\
\bottomrule
\end{tabular}
\end{table}

The label-coarsening witness in full: $C_a=\{0,1\}$, $C_b=\{2\}$, $C_c=\varnothing$, per-sample verdicts
$V=(\catch,\catch,\catch,\miss)$. Fine vocabulary (3 classes) $t_1=\{0,1\},t_2=\{2\},t_3=\{3\}$ gives blind
rates $\{0.0,0.0,1.0\}$; coarse vocabulary (1 class) gives $\{0.25\}$. All per-sample verdicts, and every
quantity the seven axioms refer to, are unchanged ($R_{\text{or}}=0.75$ throughout); only the grouped
statistics move.

\begin{table}[h]
\centering
\small
\caption{The drift algebra beside three neighbouring measurement theories. The correspondences are
analogies; we claim no physical counterpart, and A6's irreversibility is a consequence of defining $\tau$ as
column deletion.}
\label{tab:metacompare}
\begin{tabular}{@{}p{2.4cm} p{3.0cm} p{3.0cm} p{3.2cm}@{}}
\toprule
Dimension & Drift algebra & Measurement invariance & CTT / IRT \\
\midrule
Object & the caliber of an evaluator & latent variable vs.\ indicators & latent trait vs.\ item responses \\
Drift defined as & $D_\tau=\Rhat(\tau M)-\Rhat(M)$ & loadings/intercepts unequal across groups & differential item functioning \\
Reversibility & removal-type irreversible; post-processing correctable & partly recoverable by re-estimation & correctable via anchor items \\
Non-identifiability & yes (structural vs.\ real) & yes (latent scale) & yes ($\theta$ scale and origin) \\
Composability & generators compose; observed matrices do not & needs anchoring & needs equating \\
\bottomrule
\end{tabular}
\end{table}

Complexity, restated with its caveat: deciding whether a caliber change occurred is $O(|T|\cdot n)$ once the
caliber is enumerated; deciding whether the drift is structural is undecidable from paired data alone;
caliber arbitrage is NP-hard by reduction from max-coverage, with a greedy $(1-1/e)$ approximation. The
hardness statement is a reduction argument citing a classical result \citep{feige1998threshold}, not a formal
reduction proved here, and it is vacuous at $|A|=8$.

% =====================================================================
\section{Supplementary empirical tables}
\label{app:extra}

\begin{table}[h]
\centering
\small
\caption{Universality classes. ``First-order-like'' is a descriptive label for a discontinuous jump observed
at two control points, not a fitted order parameter; ``no transition'' means the response function is
monotone decreasing over the whole observed range. The exponent $\beta$ is a model-implied value (a
quadratic fit forces the mean-field value $\beta=2$), not an independent measurement.}
\label{tab:universality}
\begin{tabular}{@{}c p{2.8cm} p{4.0cm} p{2.6cm} r@{}}
\toprule
Type & Control axis & Shape & Class & $\beta$ \\
\midrule
I & continuous ($d$ degenerate assets) & continuous, smooth, concave & no transition (sub-critical) & $2.02$ \\
II & binary (E1/E2) & discontinuous jump & first-order-like & n/a \\
III & discrete (class count) & non-monotone & not a transition & n/a \\
IV & discrete (rule switch) & discontinuous jump & first-order-like & n/a \\
\bottomrule
\end{tabular}
\end{table}

\begin{table}[h]
\centering
\small
\caption{T1 on the $566$-sample split ($n=566$); ``env-gated'' abbreviates $\{$asan,ubsan,tsan$\}$, and
cw/cc abbreviate compiler-warn and cross-compile. The residual is $\le0$ in all three partitions and the
closed form matches measurement to $<10^{-6}\pp$. Row 1 attains equality, so the bound is not slack here.}
\label{tab:t1}
\begin{tabular}{@{}p{4.9cm} r r r r r@{}}
\toprule
Partition $(A\mid L\mid B)$ & $D_A$ & $D_L$ & $D_{A\cup L}$ & residual & closed form \\
\midrule
env-gated $\|$ $\{$linker$\}$ $\|$ $\{$cw,cc$\}$ & $-35.3357$ & $-0.7067$ & $-36.0424$ & $0.0000$ & $0.0000$ \\
env-gated $\|$ $\{$compiler-warn$\}$ $\|$ $\{$cc,linker$\}$ & $-35.3357$ & $-4.7703$ & $-45.0530$ & $-4.9470$ & $-4.9470$ \\
$\{$asan$\}$ $\|$ $\{$ubsan$\}$ $\|$ $\{$tsan,cw,cc,linker$\}$ & $-10.6007$ & $-7.7739$ & $-20.6714$ & $-2.2968$ & $-2.2968$ \\
\bottomrule
\end{tabular}
\end{table}

\begin{table}[h]
\centering
\small
\caption{Three-state transition matrices, $E_1\to E_2$, unaware accounting. $b=T_{\catch\to\miss}$ is the
forward McNemar discordant count; the reverse count is $c=T_{\miss\to\catch}=0$ in every frame (column
omitted). The last column reports whether the drift is silent ($\Delta_{\unk}=0$).}
\label{tab:transition}
\begin{tabular}{@{}p{3.5cm} r r r r r r r l@{}}
\toprule
Frame & $n$ & $E_1$ (c/m/u) & $E_2$ (c/m/u) & $T_{c\to c}$ & $T_{c\to m}$ & $T_{m\to c}$ & $b$ & silent \\
\midrule
A5 eval $566$ & $566$ & $340/226/0$ & $140/426/0$ & $140$ & $200$ & $0$ & $200$ & yes \\
A5 deriv $571$ & $571$ & $376/195/0$ & $158/413/0$ & $158$ & $218$ & $0$ & $218$ & yes \\
A5 all $1137$ & $1137$ & $716/421/0$ & $298/839/0$ & $298$ & $418$ & $0$ & $418$ & yes \\
cap.\ boundary $1147$ & $1147$ & $707/440/0$ & $251/896/0$ & $251$ & $456$ & $0$ & $456$ & yes \\
real-world $110$ & $110$ & $65/45/0$ & $26/84/0$ & $26$ & $39$ & $0$ & $39$ & yes \\
\bottomrule
\end{tabular}
\end{table}

\begin{table}[h]
\centering
\small
\caption{Type~I: apparent selection gain versus the number $d$ of injected zero-yield components, at
selection budget $k=4$. The greedy arm is constant at $59.55\%$ throughout---zero-yield components
contribute nothing---so the entire rise in $g_{\text{app}}$ is a denominator effect. $d\ge4$ uses injected
ghost components, not real assets, and must not be extrapolated.}
\label{tab:gapp}
\begin{tabular}{@{}rrrrr@{}}
\toprule
$d$ & assets & greedy $k{=}4$ & random-arm mean & $g_{\text{app}}$ \\
\midrule
0 & 5 & $59.55\%$ & $54.25\%$ & $5.30\pp$ \\
1 & 6 & $59.55\%$ & $48.92\%$ & $10.62\pp$ \\
2 & 7 & $59.55\%$ & $44.38\%$ & $15.16\pp$ \\
3 & 8 & $59.55\%$ & $40.50\%$ & $19.04\pp$ \\
4 & 9 & $59.55\%$ & $37.18\%$ & $22.37\pp$ \\
5 & 10 & $59.55\%$ & $34.31\%$ & $25.23\pp$ \\
6 & 11 & $59.55\%$ & $31.83\%$ & $27.71\pp$ \\
7 & 12 & $59.55\%$ & $29.67\%$ & $29.88\pp$ \\
\bottomrule
\end{tabular}
\end{table}

\begin{table}[h]
\centering
\small
\caption{Type~III: the share of classes exceeding $50\%$ blind-spot rate is \emph{not} monotone in
granularity. The finer $70$-type vocabulary contains $14$ singleton classes, which inflate the denominator
without inflating the numerator. A number of this form must be recomputed, never rescaled, when the
vocabulary changes.}
\label{tab:granularity}
\begin{tabular}{@{}p{5.6cm} r r r@{}}
\toprule
Vocabulary & classes & classes $>50\%$ blind & share \\
\midrule
\texttt{family} & $8$ & $3$ & $37.50\%$ \\
\texttt{source\_batch} & $9$ & $2$ & $22.22\%$ \\
\texttt{defect\_type} normalised & $34$ & $13$ & $38.24\%$ \\
\texttt{defect\_type} original & $70$ & $18$ & $25.71\%$ \\
\texttt{defect\_type} $\times$ \texttt{source\_batch} & $75$ & $20$ & $26.67\%$ \\
\bottomrule
\end{tabular}
\end{table}

\begin{table}[h]
\centering
\small
\caption{Type~IV: manipulability, defined as the change in the reported value produced by adding three
zero-yield components---components that add no detection capability whatsoever. OR, max and at-least-two
are immune; the per-asset mean is moved by a third; the minimum collapses. The minimum is excluded from the
health score precisely because its collapse is degenerate.}
\label{tab:manip}
\begin{tabular}{@{}p{3.6cm} r r r@{}}
\toprule
Aggregation rule & producing pool & with 3 zero-yield comps & relative change \\
\midrule
\texttt{or\_any} & $61.6391\%$ & $61.6391\%$ & $0.00\%$ \\
\texttt{max\_asset} & $35.5711\%$ & $35.5711\%$ & $0.00\%$ \\
\texttt{at\_least\_2\_assets} & $28.1604\%$ & $28.1604\%$ & $0.00\%$ \\
\texttt{mean\_per\_asset} & $17.6693\%$ & $11.7795\%$ & $\mathbf{33.33\%}$ \\
\texttt{min\_asset} & $0.8718\%$ & $0.0000\%$ & $100.00\%$ (degenerate) \\
\bottomrule
\end{tabular}
\end{table}

\begin{table}[h]
\centering
\small
\caption{Correctability by layer. ``Post-processing'' drifts leave the record $R$ intact and are repaired at
zero residual cost; ``acquisition'' drifts delete information and are not.}
\label{tab:correct}
\begin{tabular}{@{}p{6.5cm} l r r@{}}
\toprule
Drift & Layer & Correctable? & Residual \\
\midrule
accounting rule (unaware $\leftrightarrow$ aware) & post-processing & yes & $0$ \\
aggregation rule (OR $\leftrightarrow$ mean) & post-processing & yes & $0$ \\
label vocabulary & post-processing & yes & $0$ \\
structurally-$\unk$ columns & acquisition (prior known) & yes & $0$ \\
env-gated removal of asan/ubsan/tsan & acquisition & \textbf{no} & Table~\ref{tab:residual} \\
\bottomrule
\end{tabular}
\end{table}

\begin{table}[h]
\centering
\small
\caption{A four-component health score (each component $25$ points) for the instantiated apparatus. The
component values, the equal weights and the thresholds are self-set with no external calibration, so the
total should be read as a structured summary, not as a calibrated instrument.}
\label{tab:health}
\begin{tabular}{@{}lrr@{}}
\toprule
Component & raw value & score \\
\midrule
composition stability & $0.2784$ & $6.96/25$ \\
environment stability & $0.5305$ & $13.26/25$ \\
label stability & $0.8748$ & $21.87/25$ \\
aggregation stability & $0.6667$ & $16.67/25$ \\
\midrule
\textbf{total} & & $\mathbf{58.76/100}$ \\
\bottomrule
\end{tabular}
\end{table}

\paragraph{An evolution model, kept in the appendix.} Modelling the apparatus's own history as a state
vector over eight capability axes with a three-valued scale ($0$ absent, $1$ partial, $2$ complete) across
six recorded versions gives a monotone $\max$ transition with zero removals, a unique fixed point, no limit
cycles, and an undecidable chaos question (six time points is far too few for a Lyapunov estimate). The
model's role here is narrow and we keep it narrow: it is evidence that the apparatus's own evolution is
failure-driven rather than optimisation-driven, which is the SG-3 premise of this paper appearing at the
meta level. The state coordinates, the three-valued scale and the quality weights ($0.40$ detection,
$0.35$ stability, $0.25$ interpretability) are all self-set; the landscape shape should not be read as a
quantitative claim.

% =====================================================================
\section{Reproduction commands}
\label{app:repro}

Every number in this paper comes from a frozen artifact and can be recomputed. Representative one-line
commands (all read-only; no detector is invoked):

\begin{verbatim}
# T1 residual + three-state transition matrices + 692 reconciliation
python tools/compute_697_drift_algebra.py

# T5 (Chapman-Kolmogorov) + T6 (correctability residual) + health score
python tools/compute_698_drift_propagation.py

# Four-type quantification: g_app(d), per-asset sensitivity,
# granularity ladder, manipulability table
python tools/compute_698_composition_drift.py
python tools/detect_698_structural_goodhart.py --demo-queyi

# Metatheory: 4096 structures x 7 aggregations
python tools/compute_700_axiom_independence.py

# Phase-transition scan, critical point and structural early-warning signals
python tools/compute_700_phase_transition.py

# Sample complexity, KL floors, sampling-strategy simulation
python tools/compute_700_sample_complexity.py

# Cross-domain arm and migration score
python tools/compute_700_llm_drift.py
\end{verbatim}

All scripts are pure-standard-library, deterministic (no random seeds beyond fixed permutation streams) and
fail loud. The frozen matrices themselves are not regenerated by this paper.

% =====================================================================
\section{Reference verification status}
\label{app:refs}

Following the discipline of the companion paper, every bibliography entry carries a verification tag in its
\texttt{note} field. Of the entries cited here: the majority were checked online in earlier batches
(\texttt{verified-698} for the entries checked in that batch, \texttt{verified-695} for those checked
earlier); the classical books and journal articles (\texttt{classic}) have no arXiv identifier and their
metadata comes from domain knowledge; and a small number are marked \texttt{unverified-memory}, meaning
their identifiers were recalled rather than checked. \textbf{No entry marked \texttt{unverified-memory} is
used as evidence anywhere in this paper}---those are cited only as background and must be re-verified before
camera-ready. One entry (\texttt{lv2026whoevaluates}) was seen only behind a paywall and is likewise flagged
for re-verification. This paper introduces no bibliography entry that was not already present in the
repository's earlier literature batches.

\end{document}
